\documentclass[11pt]{article}

\usepackage[french, english]{babel}

\usepackage[T1]{fontenc}

\usepackage{graphicx}
\usepackage{multirow}
\usepackage{multicol}
\usepackage{xcolor}
\usepackage[shortlabels]{enumitem}

\usepackage[margin=2cm]{geometry}


\usepackage{listings}


\begin{document}

\hrule

\begin{center}
  {\Large GEI299 - Gestion de projet} \\ \vspace{5mm}
  {\LARGE\sffamily Devoir 4} \\
  Alexis Jean\\
  À remettre le 19 décembre 2025
\end{center}

\hrule

\section{\underline{Introduction}}
Dans ce document, il va y être présenté le plan de test du logiciel de simulation quantique destiné à l’analyse et à la comparaison de métaux candidats pour des agents de contraste IRM qui a été contracté par GE Healthcare. Le plan de test répondra cahier de spécification des charges fourni par le client et visera à s’assurer que le logiciel développé répond aux exigences fonctionnelles et non fonctionnelles définies.

\section{\underline{Plan de test}}
\subsection{Objectifs du plan de test}
Le plan de test a pour objectifs principaux :
\begin{itemize}
  \item Vérifier que les exigences fonctionnelles et non fonctionnelles sont respectées.
  \item Vérifier et valider les résultats de simulation produits par le logiciel.
  \item Assurer que les performances, la sécurité et la compatibilité du logiciel sont conformes aux attentes.  
  \item Vérifier l’interface utilisateur pour garantir une expérience utilisateur optimale pour les deux modes (scientifique et affaire).
  \item Trouver et régler les bugs et les problèmes potentiels avant la livraison finale du logiciel.
\end{itemize}
\subsection{Portée des tests}
Le plan de test coure l'entièreté du logiciel. Plus précisément, les tests porteront sur :
\begin{itemize}
  \item La base de données internes ainsi que le comparateur/classement des métaux candidats.
  \item Le module de simulation quantique connecté à la plateforme IBMQ.
  \item L’interface utilisateur pour les deux modes (scientifique et affaire).
  \item La fonctionnalité d'exportation des résultats.
  \item Tout ce qui n'est pas fonctionnel comme la performance, la sécurité et la compatibilité.
\end{itemize}
\newpage
\subsection{Stratégie de test}
La stratégie de tests suit une approche progressive en partant des tests unitaires jusqu'aux tests d'acceptation utilisateur. Les types de tests incluent :
\begin{itemize}
  \item \textbf{Tests unitaires :} Chaque composante individuelle du logiciel sera testé pour vérifier son bon fonctionnement.
  \item \textbf{Tests d'intégration :} Les interactions entre les différentes composantes seront testées pour s'assurer qu'elles fonctionnent ensemble correctement.
  \item \textbf{Tests système :} Le logiciel dans son ensemble sera testé pour vérifier qu'il répond aux exigences spécifiées.
  \item \textbf{Tests d'acceptation utilisateur :} Des utilisateurs finaux testeront le logiciel pour s'assurer qu'il répond à leurs besoins et attentes.
\end{itemize}
\subsection{Organisation et responsabilités}
Les différentes équipes impliquées dans le processus de test vont être répartis comme suit :
\begin{itemize}
  \item \textbf{Équipe de développement :} Responsable de l'implémentation et de la maintenance du logiciel.
  \item \textbf{Responsable scientifique :} Responsable de la validation des résultats de simulation quantique avec les références classiques et de la conformité aux exigences scientifiques.
  \item \textbf{Technicien logiciel :} Responsable de la mise en place des environnements de test, soit Windows et MacOS, et de l'exécution des tests.
  \item \textbf{Client (GE Healthcare) :} Responsable de la validation finale du logiciel et de l'acceptation des résultats.
\end{itemize}
\subsection{Environnement des tests}
Les tests seront effectués dans des environnements contrôlés qui simulent les conditions réelles d'utilisation du logiciel. Cela inclut :
\begin{itemize}
  \item Des ordinateurs sous Windows et MacOS (dans leurs dernières versions).
  \item Une connexion stable à la plateforme IBMQ pour les simulations quantiques.
  \item Un jeux de données de test représentatif des métaux candidats à analyser.
\end{itemize}
\newpage
\subsection{Description des tests}
\subsubsection{Tests fonctionnels}
\paragraph{Gestion des métaux candidats}
\begin{itemize}
  \item \textbf{But :} Vérifier l'ajout, la modification et la suppression des métaux dans la base de données.
  \item \textbf{Procédure :} L'utilisateur ajoute, modifie et supprime des métaux et vérifie que les changements sont correctement ajustés dans la base de données.
  \item \textbf{Résultat :} Les changements faits par l'utilisateur sont correctemtn appliqués dans la base de données.
\end{itemize}
\paragraph{Simulation quantique}
\begin{itemize}
  \item \textbf{But :} Valider le calcul de l'état fondamental ainsi que des propriétés électroniques des métaux candidats.
  \item \textbf{Procédure :} Éxécuter des simulations pour une liste prédéfinie de métaux et comparer les résultats avec des données de référence.
  \item \textbf{Résultat :} On trouve un écart inférieur à 2\% par rapport aux références classiques.
\end{itemize}
\paragraph{Comparateur multicritères}
\begin{itemize}
  \item \textbf{But :} Vérifier le classement automatique des métaux en fonction des critères définis.
  \item \textbf{Procédure :} Faire la comparaison de plusieurs métaux en utilisant différents critères et vérifier l'exactitude du classement.
  \item \textbf{Résultat :} Le classement des métaux correspond avec ce qui est attendu selon les critères définis.
\end{itemize}
\subsubsection{Tests de performance}
\begin{itemize}
  \item \textbf{But :} Vérifier si le logiciel respecte le temps de calcul maximum demandé par le client.
  \item \textbf{Procédure :} Mesurer le temps de calcul pour différentes simulations et s'assurer qu'il est inférieur à 2 minutes.
  \item \textbf{Résultat :} Le temps de calcul pour chaque simulation est inférieur à 2 minutes.
\end{itemize}
\subsubsection{Tests d'interface utilisateur}
\begin{itemize}
  \item \textbf{But :} Évaluer si l'interace utilisateur est clair et intuitif pour les deux modes (scientifique et affaire).
  \item \textbf{Procédure :} Faire des tâches courantes avec le logiciel dans les deux modes et recueillir les retours des utilisateurs.
  \item \textbf{Résultat :} Les tâches sont réalisés sans difficulté et sans aide extérieure.
\end{itemize}
\subsubsection{Tests de sécurité}
\begin{itemize}
  \item \textbf{But :} Vérifier la protection es accès aux données sensibles.
  \item \textbf{Procédure :} Tester l'authentification et la sécurisation de la connexion IBMQ.
  \item \textbf{Résultat :} Aucun acccès non autorisé n'est possible.
\end{itemize}
\subsubsection{Tests de compatibilité}
\begin{itemize}
  \item \textbf{But :} S'assurer que le logiciel fonctionne correctement sur les systèmes d'exploitation Windows et MacOS.
  \item \textbf{Procédure :} Faire l'ensemble des tests fonctionnels et de performance sur les deux systèmes d'exploitation.
  \item \textbf{Résultat :} Le logiciel fonctionne correctement sur les deux systèmes d'exploitation sans bugs majeurs.
\end{itemize}
\subsection{Critères d'acceptation}
Le logiciel sera considéré comme acceptable si :
\begin{itemize}
  \item Toutes les exigences fontionnelles sont respectées.
  \item Les performances et la précison respectent ce qui est attendu par le client.
  \item Aucune défaut majeur n'est trouvé lors des tests.
  \item Le client valide les test d'acceptation.
\end{itemize}
\subsection{Documentation des tests}
Tous les résultats des tests seront documentés dans des rapports tel qu'il y aura la description des tests, les résultats observés, la confirmité aux attentes ainsi que les actions qui ont été prises pour résoudre les différents problèmes.

\newpage
\section{\underline{Analyse des risques}}
L’objectif de cette section est d’identifier les risques plausibles pouvant affecter le déroulement du projet, puis d’évaluer leur importance à l’aide d’une matrice de risques qualitative afin de déterminer les actions appropriées.

\subsubsection*{Méthodologie}

Les risques sont évalués selon deux axes :
\begin{itemize}
    \item \textbf{Probabilité} : Très faible (TF), Faible (F), Moyenne (M), Élevée (E), Très élevée (TE)
    \item \textbf{Impact} : Très faible (TF), Faible (F), Moyen (M), Élevé (E), Très élevé (TE)
\end{itemize}

Le niveau de risque est déterminé par la combinaison probabilité–impact et classé selon trois catégories :
\begin{itemize}
    \item \textbf{Très faible} : risque négligeable, aucune action requise.
    \item \textbf{Faible} : risque acceptable, aucune action immédiate requise.
    \item \textbf{Modéré} : risque nécessitant un suivi.
    \item \textbf{Élevé} : risque nécessitant des actions de prévention.
    \item \textbf{Très élevé} : risque critique, actions urgentes requises.
\end{itemize}

\subsubsection*{Matrice de risques}

\begin{table}[h]
\centering
\begin{tabular}{|c|c|c|c|c|c|}
\hline
\textbf{Probabilité\textbackslash Impact} & TF & F & M & E & TE \\
\hline
TE & M & E & E & TE & TE \\
E  & F & M & E & E  & TE \\
M  & F & M & M & E  & E  \\
F  & TF & F & M & M  & E  \\
TF & TF & TF & F & F  & M  \\
\hline
\end{tabular}
\caption{Matrice de risques (5×5)}
\end{table}

\subsubsection*{Risques possibles}

\paragraph{Risque 1 – Précision insuffisante des résultats quantiques} :\\
Les résultats des simulations quantiques pourraient dépasser la tolérance d’erreur spécifiée par rapport aux références classiques.\\
Probabilité : M \quad Impact : E \quad Niveau : Élevé \\
Actions requises : validation croisée avec des cas de référence, ajustement des paramètres et tests supplémentaires.

\paragraph{Risque 2 – Retard dans l’échéancier du projet} :\\
La complexité technique et les dépendances entre tâches peuvent entraîner des retards.\\
Probabilité : E \quad Impact : M \quad Niveau : Élevé \\
Actions requises : planification avec marges de sécurité, suivi régulier de l’avancement et priorisation des tâches critiques.

\paragraph{Risque 3 – Dépassement des coûts liés à l’utilisation d’IBM Quantum} :\\
Une utilisation plus importante que prévue des ressources quantiques peut engendrer un dépassement budgétaire.\\
Probabilité : M \quad Impact : M \quad Niveau : Modéré \\
Actions requises :limitation du nombre de simulations quantiques et recours à des simulations classiques préliminaires.

\paragraph{Risque 4 – Difficultés d’utilisabilité pour les utilisateurs non scientifiques} :\\
L’interface pourrait être jugée trop complexe pour les utilisateurs du profil affaires.\\
Probabilité : F \quad Impact : M \quad Niveau : Modéré \\
Actions requises : tests d’utilisabilité et améliorations itératives de l’interface.

\paragraph{Risque 5 – Problèmes de compatibilité multiplateforme} :\\
Des différences de comportement peuvent apparaître entre les environnements Windows et MacOS.\\
Probabilité : F \quad Impact : F \quad Niveau : Faible \\
Actions requises : tests réguliers sur les deux plateformes.

\newpage
\section{\underline{Pour aller plus loin}}

\end{document}